\documentclass[11pt]{article}
\usepackage[margin=1in]{geometry}
\usepackage{amsmath,amssymb}
\usepackage{booktabs}
\usepackage[T1]{fontenc}
\usepackage[utf8]{inputenc}
\usepackage{enumitem}
\usepackage{hyperref}

\hypersetup{
  colorlinks=true,
  linkcolor=blue,
  urlcolor=blue,
  citecolor=blue
}

\setlist[itemize]{leftmargin=1.5em}
\setlist[enumerate]{leftmargin=1.7em}

\newcommand{\dt}{\Delta t}
\newcommand{\dtfine}{\Delta t_{\mathrm{fine}}}
\newcommand{\fac}{f}

\title{AthenaK Newtonian HD Local Adaptive Timestepping Implementation Note}
\author{LAT branch implementation summary}
\date{\today}

\begin{document}
\maketitle

\section{Scope}

This note documents the local adaptive timestepping (LAT) implementation in the
current \texttt{LAT} branch of this AthenaK-derived Newtonian HD code.  The
implementation is designed for the production TDE setup with LTE table EOS,
dual energy, self-gravity, AMR, and a Newtonian black-hole sink/excision
source.  It is not a general all-physics AthenaK LAT implementation.

LAT is enabled only for explicit Newtonian hydrodynamics.  The driver rejects
LAT for MHD, radiation, GR/SR coordinate systems, particles, implicit/IMEX
stages, orbital advection, shearing boxes, viscosity, conduction, generic user
timestep constraints, generic user boundary conditions, and user source terms
that are not explicitly marked LAT-safe by the problem generator.  This
restriction is intentional: every operator that changes state must be aware of
which MeshBlocks are active at the current local time.

The main implementation files are:
\begin{itemize}
  \item \path{src/driver/driver.cpp} and \path{src/driver/driver.hpp}: LAT
        window scheduling, active-mask lifetime, synchronized gravity/AMR/output
        ordering, and restart/post-AMR LAT migration.
  \item \path{src/mesh/mesh.cpp}: local-CFL metadata, factor binning,
        same-level/AMR neighbor limiting, and cached bin counts.
  \item \path{src/mesh/meshblock_pack.cpp} and
        \path{src/mesh/meshblock_pack.hpp}: active MeshBlock masks,
        boundary-provider masks, flux-correction masks, mask caches, and
        per-block LAT time intervals.
  \item \path{src/mesh/load_balance.cpp}: LAT factor estimates for initial
        load balance, default LAT GID interleaving, per-timelevel partition
        objective, and cap-aware load balancing.
  \item \path{src/mesh/build_tree.cpp},
        \path{src/mesh/mesh_refinement.cpp}, and
        \path{src/mesh/mesh_refinement.hpp}: startup/restart ordering, AMR
        candidate ordering, migration maps, and safe state transfer under
        non-Z-order GIDs.
  \item \path{src/pgen/pgen.cpp}: restart payload reads through the
        current-GID to restart-file-GID map.
  \item \path{src/hydro/hydro_tasks.cpp},
        \path{src/hydro/hydro_update.cpp},
        \path{src/hydro/hydro_newdt.cpp},
        \path{src/hydro/hydro_fluxes.cpp}, and
        \path{src/hydro/hydro_dual_energy.cpp}: active-mask-aware hydro
        kernels, delayed refluxing, and dual-energy consistency.
  \item \path{src/bvals/}: time-aware cell-centered boundary exchange,
        prolongation/restriction, and AMR/same-level flux correction.
  \item \path{src/pgen/tde_external.cpp}: LAT-safe TDE source hooks and the
        black-hole sink/excision LAT factor cap.
\end{itemize}

\section{User Controls}

The production controls are in the \texttt{<time>} and
\texttt{<mesh\_refinement>} blocks:

\begin{center}
\begin{tabular}{p{0.34\linewidth}p{0.56\linewidth}}
\toprule
Parameter & Meaning \\
\midrule
\texttt{time/hydro\_lat} & Enable HD LAT. \\
\texttt{time/hydro\_lat\_levels} & Number of binary local timestep levels.
The maximum synchronization factor is \(2^{\texttt{levels}}\). \\
\texttt{time/hydro\_lat\_same\_level} & Permit same-AMR-level neighbors to
use different LAT factors when boundary interpolation and same-level delayed
refluxing are enabled. \\
\texttt{time/hydro\_lat\_same\_level\_max\_ratio} & Maximum same-level factor
ratio after a synchronized history is available. \\
\texttt{time/hydro\_lat\_neighbor\_limiter} & Neighbor limiter scope:
\texttt{all}, \texttt{hybrid}, or \texttt{face}. \\
\texttt{time/hydro\_lat\_min\_bin\_count} & Minimum useful bin population.
\(-1\) means the default \(4N_{\mathrm{rank}}\); \(0\) disables sparse-bin
pruning. \\
\texttt{time/hydro\_lat\_gid\_reorder} & Enable LAT-optimized GID ordering.
Default: \texttt{true} when LAT is enabled. \\
\texttt{time/hydro\_lat\_post\_amr\_rebalance} & Allow the driver to perform
a migration-only LAT rebalance after AMR when needed.  The default follows
\texttt{hydro\_lat\_gid\_reorder}. \\
\texttt{time/hydro\_lat\_diagnostics} & Print additional LAT bin and factor
diagnostics. \\
\texttt{mesh\_refinement/max\_nmb\_per\_rank} & Hard GPU memory cap on
MeshBlocks per rank, honored by LAT startup, restart, load balance, and AMR.
For the current H2 production input this is set to 600. \\
\bottomrule
\end{tabular}
\end{center}

When \texttt{hydro\_lat=true}, the driver internally uses hydro subcycling with
\[
  \texttt{hydro\_subcycle\_factor} = 2^{\texttt{hydro\_lat\_levels}} .
\]
This is an upper bound.  Each LAT window chooses a possibly smaller
synchronization factor from the available metadata, the remaining time to
\texttt{tlim}, and the self-gravity cadence.

\section{Basic Time Integration Algorithm}

At a synchronized state all MeshBlocks are at the same physical time \(t^n\).
The code computes the ordinary global hydro timestep \(\dtfine\), which is the
minimum local-CFL hydro timestep over all MeshBlocks.  It then assigns each
MeshBlock a power-of-two local factor
\[
  \fac_i \in \{1,2,4,\ldots,F\},
\]
where \(F\) is the active synchronization factor for the next LAT window.
MeshBlock \(i\) advances with
\[
  \dt_i = \fac_i \dtfine .
\]

The driver then advances \(F\) fine ticks.  At fine tick phase \(q\), a block
with factor \(\fac\) is due if
\[
  q \bmod \fac = 0 .
\]
Within each fine tick, due bins are visited from the largest factor down to
\(\fac=1\).  This coarsest-first sweep makes slow blocks provide boundary data
before faster bins need it at the same fine-tick time.

At the end of every fine tick the mesh time and cycle advance by
\(\dtfine\), not by the largest bin timestep.  At the end of the LAT window all
blocks are synchronized again.  Only synchronized states are allowed to trigger
AMR, output, after-cycle callbacks, local-CFL metadata rebuilds, or multigrid
self-gravity source loading.

\section{LAT Factor Metadata}

\path{Hydro::NewTimeStep} stores local hydro timestep estimates in
\texttt{dtnew\_eachmb}.  At synchronized states,
\path{Mesh::UpdateHydroLATMetadata} gathers these estimates into
\texttt{hydro\_lat\_dt\_eachmb} and builds
\texttt{hydro\_lat\_factor\_eachmb}.  The raw factor is the largest
power-of-two value satisfying
\[
  \fac_i \dtfine \leq \dt_{i,\mathrm{local}},
\]
clamped by the configured maximum factor.

The raw factors are then made safe by several reductions:
\begin{enumerate}
  \item A problem-specific cap hook may lower factors when a problem generator
        installs one.  The TDE problem generator currently does not install a
        LAT factor cap.
  \item If \texttt{hydro\_lat\_same\_level=false}, all blocks on an AMR level
        collapse to the minimum factor on that level.
  \item Sparse bins are folded downward according to
        \texttt{hydro\_lat\_min\_bin\_count}.  The default is
        \(4N_{\mathrm{rank}}\).
  \item The neighbor limiter enforces a conservative staircase across AMR and
        same-level neighbors.  The limiter is iterative and MPI-reduced: local
        changes are detected with a small allreduce first, and the full factor
        array is reduced only while changes still exist.
\end{enumerate}

The final factor table is versioned with
\texttt{hydro\_lat\_metadata\_version}.  MeshBlockPack active-mask caches are
invalidated whenever metadata or topology changes.  The mesh also stores
\texttt{hydro\_lat\_bin\_count}, so common per-tick queries such as ``does bin
\(\fac\) exist'' are \(O(1)\) instead of repeatedly scanning
\texttt{nmb\_total}.

\section{Active Masks and Communication Masks}

LAT is implemented by masking MeshBlocks inside one ordinary MeshBlockPack.
The main arrays live in \path{MeshBlockPack}:
\begin{itemize}
  \item \texttt{lat\_active\_mb} and \texttt{lat\_active\_indices}: blocks that
        execute state-changing hydro work in the current bin.
  \item \texttt{lat\_boundary\_send\_mb} and
        \texttt{lat\_boundary\_send\_indices}: inactive blocks that still need
        to provide ghost-zone data to active neighbors.
  \item \texttt{lat\_send\_nghbr}: per-neighbor boundary send mask.
  \item \texttt{lat\_flux\_accum\_nghbr},
        \texttt{lat\_flux\_send\_nghbr}, and
        \texttt{lat\_flux\_recv\_nghbr}: face-neighbor masks for delayed
        refluxing.
  \item \texttt{lat\_time\_start} and \texttt{lat\_time\_end}: the local time
        interval represented by each block's current conserved state.
\end{itemize}

\path{SetActiveMeshBlocksByLATFactor} installs a single factor bin.
\path{SetActiveMeshBlocksByLATDueFactors} installs the union of all bins due at
a fine-tick phase.  The union mask is used for the pre-bin boundary refresh at
the tick start.  Both routines compute and cache the boundary-need flags
\texttt{lat\_cached\_needs\_restrict},
\texttt{lat\_cached\_needs\_prolongate},
\texttt{lat\_cached\_needs\_physical\_bc}, and
\texttt{lat\_cached\_needs\_neighbor\_recv}.  Therefore
\path{RefreshHydroLATBoundaries} does not have to rescan all neighbors for
every refresh.

It is important that ``active'' and ``communication needed'' are different.
A rank may have no active MeshBlocks in a bin but still own boundary providers
or flux receivers.  In that case the relevant task lists still run on that
rank so MPI send/receive ordering remains consistent.

\section{Boundary Exchange and Time Interpolation}

Cell-centered boundary exchange is time-aware under LAT.  The driver sets
\texttt{hydro\_lat\_exchange\_time} before a refresh.  Boundary packing and
unpacking can use the block's start state, end state, and requested target time
to provide ghost data at the correct local time.  Same-level mixed-factor
boundaries use dense RK2 interpolation between the synchronized start state
and the block's current/end state.

The driver refresh sequence is:
\begin{enumerate}
  \item Optional problem hydro state fixup at the target time.
  \item Restriction if an active or provider block has a coarser neighbor that
        needs restricted data.
  \item The usual \texttt{InitRecv}, \texttt{SendU}, \texttt{RecvU} boundary
        exchange and shear exchange.
  \item Physical BCs, prolongation, and conserved-to-primitive conversion only
        when the cached mask says the operation is needed.
\end{enumerate}

Non-fine bins may skip their final boundary exchange when the next fine tick
will perform a union refresh.  This avoids redundant communication without
leaving synchronized states stale: the driver marks a boundary refresh as
needed after each fine tick and performs the refresh at the next synchronized
state before gravity, AMR, output, or timestep recomputation.

\section{Delayed Flux Correction}

Finite-volume conservation across mixed-timestep faces is handled by delayed
refluxing.  During a LAT bin,
\path{Hydro::AccumulateLATCoarseFluxes} accumulates the slower/coarser side's
time-integrated face flux using the final RK flux weight and the local
timestep.  Faster/finer-side fluxes are sent through the existing AMR
flux-correction buffers.  When a receiver reaches the corresponding
synchronization time, \path{Hydro::ApplyLATFluxCorrection} applies the flux
difference.

This mechanism is used for both ordinary AMR coarse/fine faces and same-level
mixed-factor faces when \texttt{hydro\_lat\_same\_level=true}.  Edge and corner
neighbors participate in boundary exchange, but not flux correction, because
finite-volume fluxes pass only through faces.

The reflux correction is split into independent x1, x2, and x3 directional
kernels with fences between directions.  This avoids races in corner/edge
cells when multiple face corrections touch the same conserved cell.  For
dual-energy hydro, the auxiliary internal-energy/volume-fraction field receives
a matching delayed correction.  For self-gravity total-energy runs, the reflux
path includes a gravitational energy consistency correction using the current
potential.

\section{Hydro, EOS, and Source-Term Coverage}

All state-changing hydro kernels used by this LAT mode must respect
\texttt{lat\_active\_mask\_enabled}.  The current implementation covers RK
updates, flux computation, dual-energy synchronization, EOS
conserved-to-primitive conversion, physical hydro boundary work, AMR
prolongation/restriction, FOFC-related hydro paths, excision cleanup, and the
TDE source term.

\path{Hydro::NewTimeStep} is not run inside masked LAT bins.  Local-CFL
estimates are meaningful only at synchronized states after boundary and reflux
corrections are complete.

Generic user source terms are rejected unless the problem generator sets
\texttt{user\_srcs\_lat\_safe=true}.  The TDE problem generator sets this only
when \texttt{problem/use\_translating\_frame=false}.  The translating frame
uses global mesh-state-dependent acceleration and is therefore not LAT-safe.
The active TDE gravity source loops over \texttt{lat\_active\_indices} when LAT
is active and evaluates BH/frame data at the driver-provided stage/source time.

\section{Self-Gravity}

Self-gravity is synchronized with LAT.  A LAT window reuses one synchronized
potential; it does not run Poisson solves inside local bins, because doing so
would load density from blocks at mixed local times.  Before a LAT window, the
driver refreshes the potential if it is invalid or due at the window start.
It also shortens the proposed LAT window if the configured gravity cadence
would require a solve before the proposed synchronization point.

At a synchronized LAT boundary the driver clears masks, refreshes hydro
boundaries if delayed reflux changed any state, and then runs gravity if due.
AMR and output see only synchronized hydro and gravity state.

\section{LAT-Aware GID Ordering}

The most important production load-balance change is the default LAT GID
ordering.  Standard AthenaK assigns GIDs by Z-order traversal, and each rank
owns a contiguous GID interval.  That contiguous interval assumption is used in
many MPI and local-index calculations, so this branch does not assign arbitrary
non-contiguous GIDs to ranks.

Instead, when \texttt{hydro\_lat=true} and
\texttt{hydro\_lat\_gid\_reorder=true} (the default), the code changes the GID
ordering itself.  \path{BuildHydroLATGIDMap} groups candidate MeshBlocks by
LAT factor, sorts each factor bin by Z-order for locality, and then emits a
proportional interleaving of factor bins.  \path{ApplyHydroLATGIDMapToArrays}
applies that permutation to logical locations, costs, factor lists, optional
restart maps, and the MeshBlockTree leaf GIDs.

This keeps the invariant
\[
  \mathrm{local\ lid} = \mathrm{gid} - \mathrm{gids\_eachrank[rank]}
\]
valid everywhere.  It also distributes fine-timestep factor-1 blocks across
all ranks, instead of letting a compact BH/stream region occupy only one or two
contiguous Z-order rank intervals.  Existing output readers continue to work
because output carries each MeshBlock's logical metadata and data; the internal
GID order is not the physical plotting coordinate system.

\section{LAT Load Balance Objective}

After GID ordering, \path{Mesh::LoadBalance} still produces contiguous GID
cuts.  When factor metadata is available, the objective is per-timelevel rather
than purely scalar.  It builds prefix tables for:
\begin{itemize}
  \item each factor bin count,
  \item scalar cost,
  \item work due at every tick phase in the LAT synchronization window.
\end{itemize}

The optimizer evaluates partitions by maximum per-tick work, sum of per-tick
rank maxima, idle rank slots, factor-bin idle slots, maximum per-rank block
count, block-count imbalance, and scalar cost.  It tries several seeds:
ordinary scalar cost, active work, equal blocks, factor-weighted costs,
peak-tick weights, bin-balanced cuts, and critical-tick cuts.  It then applies
single-cut and adjacent-pair local search while enforcing
\texttt{max\_nmb\_per\_rank}.

If LAT metadata is unavailable, startup/restart load balance estimates factors
conservatively from existing metadata, AMR containment, or AMR level.  If even
that is unavailable, it falls back to block-count balance so the memory cap
remains a true allocation bound.

\section{AMR and State Migration}

AMR is allowed only at synchronized LAT states.  The driver also suppresses a
terminal AMR pass after the last accepted hydro update, since there is no
following step to use the adapted mesh.

\path{MeshRefinement::RedistAndRefineMeshBlocks} has been made safe for
LAT-reordered GIDs.  For pure rebalance, it can request a new LAT GID map
directly from the current logical-location list.  For real AMR, it first lets
the tree produce the usual candidate mesh and new-to-old map, then applies the
LAT interleaving to both the candidate logical-location list and the
\texttt{newtoold} migration map before \texttt{oldtonew}, load balance, MPI
migration, and local in-place copies are built.

The AMR state-transfer routines were also adjusted for non-Z-order GIDs:
\begin{itemize}
  \item \texttt{oldtonew} is built by logical containment rather than by
        assuming adjacent old/new GID arithmetic.
  \item Refine/derefine child lookup uses explicit child maps rather than
        relying on consecutive child GIDs.
  \item Local copy/rebalance routines handle arbitrary same-rank source/dest
        permutations with temporary storage for cycles.
  \item After topology changes, the MeshBlockTree leaf GIDs are reset from the
        final LAT-ordered logical-location array.
\end{itemize}

Hydro storage resize was changed to reduce transient GPU memory spikes during
AMR.  Persistent state arrays may retain capacity, while scratch and
time-evolving arrays use \texttt{Kokkos::realloc} to avoid holding both old and
new scratch allocations longer than needed.

\section{Restart Semantics}

LAT restarts require a single global restart file.  Per-rank restart files are
rejected when LAT is enabled because arbitrary GID reordering cannot safely
read rank-local payload files without a global data map.

Restart handling has two distinct maps:
\begin{enumerate}
  \item The restart file stores payloads by the GID order used when the file
        was written.
  \item The live mesh may immediately reorder GIDs for LAT load balance.
\end{enumerate}
Therefore \path{BuildTreeFromRestart} constructs
\texttt{restart\_gid\_eachmb}, a current-GID to restart-file-GID map.  If the
startup LAT GID interleaving composes another permutation, that restart map is
composed as well.  \path{ProblemGenerator::ProblemGenerator} then reads every
cell-centered and face-centered restart payload from
\[
  \texttt{headeroffset}
  + \texttt{restart\_file\_gid}\times\texttt{data\_size}
  + \texttt{field\_offset}.
\]
The mapped path uses independent \texttt{MPI\_File\_read\_at} calls, so ranks
do not have to issue collective reads in identical GID order.  After all
payloads are read, the restart map is deleted.

On restart, the driver computes a fresh timestep and LAT metadata.  If the
mesh is already LAT-reordered, it keeps that ordering.  If not, it may perform
a migration-only \path{RedistAndRefineMeshBlocks(pin,0,0)} pass to move the
live state into the LAT-optimized ordering before normal evolution.

\section{TDE-Specific Sink and Source Safety}

The TDE problem generator no longer installs a BH-specific LAT factor cap.
Sink-adjacent blocks are handled by the normal LAT factor pipeline and neighbor
limiter.

The BH gravity and floor-shell cleanup source loops use the LAT active-index
list when LAT is active.  The source time comes from the driver's stage/source
time hook.  Translating-frame runs are marked LAT-unsafe, because the frame
acceleration depends on global gas state and would otherwise be evaluated
inconsistently across local times.

\section{Performance Notes}

The current branch includes several performance-oriented changes that are
important to production scaling:
\begin{itemize}
  \item Factor-bin counts are cached in \texttt{hydro\_lat\_bin\_count},
        avoiding repeated full-array scans during each LAT tick.
  \item Active-mask and due-factor masks are cached by metadata version,
        maximum factor, tick phase, and neighbor-table size.
  \item Boundary refresh gating uses cached need flags computed during mask
        construction rather than rescanning neighbors on every refresh.
  \item Neighbor limiter convergence first reduces a small changed flag before
        reducing the full factor array.
  \item The LAT load balancer prints per-bin and per-tick efficiency diagnostics
        when active, making fine-tick rank imbalance directly visible.
\end{itemize}

The dominant cost that remains is communication and task-list launch overhead
for frequent boundary refreshes at large synchronization factors.  The design
intentionally favors correctness and synchronized AMR/gravity behavior over
fully asynchronous per-rank progress.

\section{Current Production Configuration}

The H2 production input currently uses:
\begin{itemize}
  \item \texttt{hydro\_lat=true},
  \item \texttt{hydro\_lat\_levels=7}, giving a configured maximum factor
        \(2^7=128\),
  \item \texttt{hydro\_lat\_same\_level=true},
  \item \texttt{hydro\_lat\_same\_level\_max\_ratio=8},
  \item \texttt{hydro\_lat\_neighbor\_limiter=face},
  \item default \texttt{hydro\_lat\_gid\_reorder=true},
  \item \texttt{max\_nmb\_per\_rank=600}.
\end{itemize}

Observed restart/load-balance diagnostics for this configuration have shown
LAT GID interleaving across eight factor bins with worst-tick efficiency close
to unity in the initial restarted mesh.  The exact factor distribution is
runtime-dependent because it follows the local-CFL table and AMR topology.

\section{Limitations and Invariants}

This implementation should not be treated as a full H-AMR replacement.  Its
current limitations are:
\begin{itemize}
  \item Newtonian HD only.
  \item RK1/RK2 are the validated integrators for the delayed-reflux path.
  \item AMR, gravity, output, callbacks, and timestep metadata are synchronized
        at LAT window boundaries, not run asynchronously in local bins.
  \item Same-level temporal interpolation is dense RK2 between start and end
        conserved states; very large same-level ratios remain a physics
        validation risk.
  \item Per-rank restart files are not supported with LAT.
\end{itemize}

Future LAT changes should preserve these invariants:
\begin{enumerate}
  \item Never run AMR, output, self-gravity source loading, or after-cycle
        callbacks on a partially synchronized LAT mask.
  \item Recompute local-CFL LAT metadata only at synchronized states.
  \item Keep active evolution masks separate from boundary-provider and
        flux-receiver masks.
  \item Keep MeshBlock GID ownership contiguous per rank; use GID ordering
        permutations rather than arbitrary non-contiguous rank assignments.
  \item Compose restart payload maps whenever startup GIDs are permuted.
  \item Keep the AMR \texttt{newtoold}/\texttt{oldtonew} maps consistent with
        any LAT candidate ordering before MPI migration or local copies.
  \item Preserve the hard \texttt{max\_nmb\_per\_rank} cap and
        \texttt{NUM\_BITS\_LID} tag limit.
  \item Do not mark a user source LAT-safe until every state-changing loop is
        active-mask-aware and all time-dependent data use the driver-provided
        stage/source time.
  \item If same-level LAT behavior is extended, update boundary interpolation
        and delayed refluxing together.
\end{enumerate}

\end{document}
